\subsection{The unramified corollary: the generic lifting condition}

We now close the unramified corollary independently of the remaining
actual-versus-pro-modular issue in the locally quasi-finite theorem.
Throughout this subsection, condition 8 is tested on schemes locally of
finite type over its noetherian test base, and NR(b) has the unique
associated prime hypothesis of Proposition~\ref{A:NR-density}.

\begin{lemma}[The Artinian base case]\label{A:artinian-base}
\status{Proved under conditions 1--5}
Let $A$ be a local Artinian $S$-algebra. The restriction $F_A$ is
represented by a disjoint union of spectra of finite local $A$-algebras.
If the Artinian version of NR(a) holds, these algebras are unramified
over $A$.
\end{lemma}
\begin{proof}
Use the universal prorepresentability in condition 5, with the
minimal fields of definition included in its universal formulation.
It gives one local representing algebra $B_i$ for each point of the
Artinian-base functor; condition 5 says that $B_i$ is finite over $A$.
Thus $B_i$ is Artinian, and its prorepresenting element is an actual
element $\xi_i\in F(B_i)$. The morphism
\[
 h_{\coprod_i\Spec B_i}\longrightarrow F_A
\]
is bijective on local Artinian $A$-algebras, including extensions of
the residue field, by that universal property. Both sides satisfy
conditions 1--4, so Lemma~\ref{A:comparison} makes it an isomorphism
on all $A$-schemes. The Artinian infinitesimal injectivity assertion
in NR(a) applies to the represented functor and hence to each $B_i$.
For a finite algebra, formal unramifiedness is equivalent to
unramifiedness. No algebraization or excellence hypothesis enters
this argument. This is also Murre's \S2, Step 1, with finiteness
supplied by condition 5 rather than by unramifiedness.
\end{proof}

\begin{proposition}[Condition 8 implies Murre's (F8)]\label{A:Murre-F8}
\status{Proved}
Assume conditions 1--5 and 8 and the Artinian unramifiedness condition
NR(a). Then $F$ satisfies condition (F8) of Murre's Expos\'e 294,
\S1. In fact, the same generic nonextension conclusion holds without
the additional restriction $NI=0$ in (F8).
\end{proposition}
\begin{proof}
We state all the data of the test. Let $A$ be noetherian with
irreducible spectrum, let $\mathfrak p$ be its unique minimal prime,
let $I\subset A$ be a nilpotent ideal, and fix
$\eta_0\in F(A/I)$. Suppose its restriction to
$\Spec(A_{\mathfrak p}/I_{\mathfrak p})$ does not extend to any
strictly larger closed subscheme of $\Spec A_{\mathfrak p}$.
We must find a nonempty open $W\subset\Spec A$ such that the
analogous nonextension statement holds on every open $W_1\subset W$.

The ring $A_{\mathfrak p}$ is Artinian. By
Lemma~\ref{A:artinian-base}, $F_{A_{\mathfrak p}}$ is represented by
a disjoint union of finite local unramified schemes. The point
selected by $\eta_0$ has residue field $\kappa(\mathfrak p)$: a
point defined over this field in a finite unramified scheme must be
rational. Its local component therefore has the form
\[
 X_0=\Spec(A_{\mathfrak p}/J)
 \quad\text{with}\quad J\subseteq I_{\mathfrak p}.
\]
Indeed, a finite local unramified algebra over an Artinian local
ring with identical residue field is a quotient: its maximal ideal
is generated by the base maximal ideal, and successive applications
of Nakayama to this nilpotent filtration show that the image of $1$
generates it as a module over the base. The structural ring map is
therefore surjective. The inclusion $J\subseteq I_{\mathfrak p}$
follows from the given section.

The universal element on $X_0$ is an extension of $\eta_0$ to
that intermediate closed subscheme. The assumed generic
nonextension forces $J=I_{\mathfrak p}$. Consequently the generic
local scheme of $Z=\Spec(A/I)$ is an actual local representing
scheme at this point. This is local modularity for $F_A$, not just
pro-modularity: every local $A$-test whose closed point lies over
$\mathfrak p$ factors through $A_{\mathfrak p}$, and the displayed
component represents the relevant subfunctor on all such tests.
Its rational residue point also supplies the field-of-definition
condition in cote 131, scan 35.

Apply condition 8 to the irreducible noetherian base $\Spec A$ and
its finite closed subscheme $Z$. It gives a nonempty open of $Z$
on which $\eta_0$ is locally modular. Since $I$ is nilpotent,
$|Z|=|\Spec A|$, so this is an open $W\subset\Spec A$.
Let $W_1\subset W$ be open. Suppose that an element on an
intermediate closed subscheme $Z_1$ of $W_1$, containing
$Z\cap W_1$, restricts to $\eta_0$. At any point $w\in W_1$
write the ideal of $Z_1$ locally as $J_w\subseteq I_w$.
Local modularity gives a local $A$-morphism
\[
 \Spec(A_w/J_w)\longrightarrow\Spec(A_w/I_w)
\]
inducing the given element, since their closed-point values agree.
Such a morphism over $A$ forces $I_w\subseteq J_w$.
Thus $J_w=I_w$ at every point, and the two closed subschemes are
equal. This proves the required nonextension on every $W_1$.
The proof did not use $NI=0$.
\end{proof}

\begin{corollary}[Cote 131, Corollaire 2, with corrected NR(b)]
\label{A:unramified-corollary}
\status{Known criterion, with the manuscript-to-source implication proved}
For $S$ locally noetherian, conditions 1--6 and 8 together with
NR(a) and corrected NR(b) are necessary and sufficient for $F$ to
be represented by a separated unramified $S$-scheme locally of
finite type. The Artinian version of NR(a) suffices.
\end{corollary}
\begin{proof}
We apply Murre, Expos\'e 294, \S1, Theorem 1
\cite{Murre1965}, whose hypotheses we now match explicitly.
Conditions 1 and 2 give (F1), condition 3 gives (F2), and
condition 4 gives (F3). Universal Artinian prorepresentability
in condition 5 implies the special left-exactness test (F4).
Lemma~\ref{A:artinian-tests} promotes NR(a) to formal
unramifiedness on all nilpotent thickenings, giving (F5).
Valuative uniqueness is (F6). Corrected NR(b) is precisely (F7):
the complete local one-dimensional test has only one associated
prime and $NI=0$. Proposition~\ref{A:Murre-F8} gives (F8).
Thus every hypothesis of Murre's theorem holds, and its conclusion
is exactly the asserted separated unramified scheme, including
local finite type.

Conversely, a representing separated unramified scheme has
conditions 1--6; its finite infinitesimal local algebras give the
finite form of condition 5. Proposition~\ref{A:8-literal} gives
condition 8 in the finite-type formulation. Formal
unramifiedness gives NR(a), and
Proposition~\ref{A:NR-density} gives corrected NR(b).
The correction is essential when the manuscript's word is read
as ``irreducible''. If it were read as ``integral'', the displayed
nilpotent ideal test would be vacuous and would not replace (F7);
we make no claim that this weaker reading suffices.
\end{proof}

This closes Corollaire 2 without assuming the unproved local-model
premise of the locally quasi-finite theorem. It also identifies
exactly how the manuscript's generic modularity condition supplies
the published generic nonextension axiom. Murre's Note after
Corollary 2 asserts the possibility of a locally quasi-finite
analogue but does not supply its axioms or a proof; we do not use
that Note as a proof of the remaining main theorem.
